\documentclass[12pt, a4paper]{report}

% ─── Paquetes ───────────────────────────────────────────────────────────────
\usepackage[utf8]{inputenc}
\usepackage[T1]{fontenc}
\usepackage[spanish]{babel}
\usepackage{geometry}
\usepackage{graphicx}
\usepackage{hyperref}
\usepackage{booktabs}
\usepackage{listings}
\usepackage{xcolor}
\usepackage{enumitem}
\usepackage{float}
\usepackage{amsmath}
\usepackage{fancyhdr}
\usepackage{titlesec}
\usepackage{parskip}
\usepackage{microtype}
\usepackage{caption}
\usepackage{subcaption}
\usepackage{longtable}
\usepackage{array}
\usepackage{multirow}
\usepackage{lmodern}

% ─── Geometría ──────────────────────────────────────────────────────────────
\geometry{
  top=2.5cm,
  bottom=2.5cm,
  left=3cm,
  right=2.5cm
}

% ─── Colores para código ────────────────────────────────────────────────────
\definecolor{codegray}{rgb}{0.95,0.95,0.95}
\definecolor{codegreen}{rgb}{0.0,0.5,0.0}
\definecolor{codeblue}{rgb}{0.0,0.0,0.8}
\definecolor{codepurple}{rgb}{0.5,0.0,0.5}

\lstdefinestyle{pythonstyle}{
  backgroundcolor=\color{codegray},
  commentstyle=\color{codegreen},
  keywordstyle=\color{codeblue}\bfseries,
  stringstyle=\color{codepurple},
  basicstyle=\ttfamily\footnotesize,
  breakatwhitespace=false,
  breaklines=true,
  captionpos=b,
  keepspaces=true,
  numbers=left,
  numbersep=5pt,
  numberstyle=\tiny\color{gray},
  showspaces=false,
  showstringspaces=false,
  showtabs=false,
  tabsize=2,
  language=Python
}
\lstset{style=pythonstyle}

% ─── Encabezado y pie ───────────────────────────────────────────────────────
\pagestyle{fancy}
\fancyhf{}
\fancyhead[L]{\small Sistema RAG para la Gestión de Conocimiento Técnico}
\fancyhead[R]{\small FIUBA --- 2026}
\fancyfoot[C]{\thepage}
\renewcommand{\headrulewidth}{0.4pt}

% ─── Hipervínculos ──────────────────────────────────────────────────────────
\hypersetup{
  colorlinks=true,
  linkcolor=black,
  urlcolor=blue,
  citecolor=black
}

% ════════════════════════════════════════════════════════════════════════════
\begin{document}

% ─── CARÁTULA ───────────────────────────────────────────────────────────────
\begin{titlepage}
  \centering
  \vspace*{0.5cm}

  \includegraphics[width=0.35\textwidth]{uba_logo.png}\\[1.2cm]

  {\scshape\large Universidad de Buenos Aires\\[0.2cm]
  Facultad de Ingeniería\\[0.2cm]
  Año 2026\par}

  \vspace{2.5cm}

  {\scshape\Large Trabajo Profesional --- Sistema RAG para la\\[0.3cm]
  Gestión de Conocimiento Técnico\par}

  \vspace{0.5cm}
  {\large\itshape Recuperación inteligente de información histórica de cuencas maduras\par}

  \vspace{3cm}

  \begin{flushleft}
    \textbf{INTEGRANTES:}\\[0.4cm]
    \begin{tabular}{lll}
      \#108456 & Finci, Dante Alejandro       & \href{mailto:dfinci@fi.uba.ar}{dfinci@fi.uba.ar} \\
      \#110883 & Lanfranco, Tomás             & \href{mailto:tomilanfranco04@gmail.com}{tomilanfranco04@gmail.com} \\
      \#102815 & Rodriguez, Franco Ezequiel   & \href{mailto:frrodriguez@fi.uba.ar}{frrodriguez@fi.uba.ar} \\
      \#109525 & Urbano, Sol Guadalupe        & \href{mailto:surbano@fi.uba.ar}{surbano@fi.uba.ar} \\
    \end{tabular}
  \end{flushleft}

  \vspace{1.5cm}

  \begin{flushleft}
    \textbf{Tutor:} Dr.\ Hernán Daniel Merlino --- \href{mailto:hmerlino@fi.uba.ar}{hmerlino@fi.uba.ar}\\[0.3cm]
    \textbf{Co-tutora:} Dra.\ Gabriela Beatriz Savioli
  \end{flushleft}

\end{titlepage}

% ─── ÍNDICE ─────────────────────────────────────────────────────────────────
\tableofcontents
\newpage

% ════════════════════════════════════════════════════════════════════════════
\chapter*{1. Resumen}
\addcontentsline{toc}{chapter}{1. Resumen}

% TODO: Completar con el resumen final del trabajo.
\textit{[Pendiente --- completar al finalizar el proyecto.]}

% ════════════════════════════════════════════════════════════════════════════
\chapter*{2. Abstract}
\addcontentsline{toc}{chapter}{2. Abstract}

% TODO: Completar con el abstract en inglés.
\textit{[Pendiente --- completar al finalizar el proyecto.]}

% ════════════════════════════════════════════════════════════════════════════
\chapter*{3. Palabras clave / Keywords}
\addcontentsline{toc}{chapter}{3. Palabras clave / Keywords}

\textbf{Palabras clave:} RAG, gestión de conocimiento técnico, cuencas maduras, LLMs,
base de datos vectorial, ETL, OCR, pozos análogos, recuperación semántica.

\textbf{Keywords:} RAG, technical knowledge management, mature basins, LLMs,
vector database, ETL, OCR, analog wells, semantic retrieval.

% ════════════════════════════════════════════════════════════════════════════
\chapter*{4. Agradecimientos}
\addcontentsline{toc}{chapter}{4. Agradecimientos}

% TODO: Completar con agradecimientos.
\textit{[Pendiente.]}

% ════════════════════════════════════════════════════════════════════════════
\chapter*{5. Introducción}
\addcontentsline{toc}{chapter}{5. Introducción}

En la industria del petróleo y gas, la explotación de cuencas maduras genera a lo largo del
tiempo grandes volúmenes de información técnica: intervenciones realizadas, decisiones tomadas
y resultados obtenidos quedan registrados en informes, diagramas de completación y registros
operativos.

Esta documentación representa un activo de conocimiento valioso, ya que refleja décadas de
experiencia concreta sobre el comportamiento de los yacimientos. Sin embargo, se encuentra
dispersa y en formatos no estructurados, lo que dificulta su acceso y aprovechamiento,
limitando la reutilización del conocimiento experto.

El objetivo de este trabajo es mejorar el acceso y la recuperación de esa información. Para ello
se propone el desarrollo de un sistema basado en una arquitectura de Generación Aumentada
por Recuperación (RAG), apoyado en modelos de lenguaje de gran escala (LLMs). Asimismo,
se contempla el uso de una base de datos vectorial junto con procesos de ingesta de datos (ETL)
que incorporan OCR para el tratamiento de documentos históricos.

% ════════════════════════════════════════════════════════════════════════════
\chapter*{6. Estado del Arte}
\addcontentsline{toc}{chapter}{6. Estado del Arte}

\section*{6.1 Modelos de Lenguaje de Gran Escala (LLMs)}

% TODO: Completar con investigación de Fran sobre LLMs.
\textit{[Ver sección de Investigación --- Franco Rodriguez.]}

\section*{6.2 Arquitecturas RAG}

% TODO: Completar.
\textit{[Pendiente.]}

\section*{6.3 Bases de Datos Vectoriales}

% TODO: Completar con investigación de Dante sobre bases de datos vectoriales.
\textit{[Ver sección de Investigación --- Dante Finci.]}

\section*{6.4 Embeddings y Representación Semántica}

% TODO: Completar con investigación de Sol sobre embeddings.
\textit{[Ver sección de Investigación --- Sol Urbano.]}

\section*{6.5 OCR y Procesamiento de Documentos}

% TODO: Completar con investigación de Tomi sobre ETL y OCR.
\textit{[Ver sección de Investigación --- Tomás Lanfranco.]}

\section*{6.6 Aplicaciones en la Industria del Petróleo y Gas}

% TODO: Completar con referencias a trabajos similares en el dominio.
\textit{[Pendiente.]}

% ════════════════════════════════════════════════════════════════════════════
\chapter*{7. Problema Detectado y/o Faltante}
\addcontentsline{toc}{chapter}{7. Problema Detectado y/o Faltante}

En el contexto de cuencas maduras, las decisiones técnicas tomadas a lo largo de décadas
quedan registradas en grandes volúmenes de documentación: informes de intervenciones,
diagramas de completación y registros operativos. Esta información representa un activo
importante, ya que refleja qué funcionó y qué no en situaciones concretas a lo largo del tiempo.

El problema es que dicha documentación no está centralizada ni estructurada. Consultar
antecedentes de un pozo o buscar experiencias previas en situaciones similares requiere revisar
manualmente archivos dispersos. En la práctica, esto lleva a que ese conocimiento acumulado no
se aproveche de forma eficiente, aumentando el riesgo de repetir errores ya documentados y
dificultando la transferencia de experiencia entre profesionales del sector.

% ════════════════════════════════════════════════════════════════════════════
\chapter*{8. Solución Implementada}
\addcontentsline{toc}{chapter}{8. Solución Implementada}

\section*{8.1 Arquitectura Diseñada}

% TODO: Completar con el diagrama de arquitectura final y descripción.
\textit{[Pendiente --- incluir diagrama de arquitectura.]}

\section*{8.2 Pipeline ETL y OCR}

% TODO: Completar con la implementación del pipeline.
\textit{[Pendiente.]}

\section*{8.3 Generación de Embeddings}

% TODO: Completar con el modelo de embeddings seleccionado y justificación.
\textit{[Pendiente.]}

\section*{8.4 Base de Datos Vectorial}

% TODO: Completar con la implementación de la base de datos vectorial.
\textit{[Pendiente.]}

\section*{8.5 Integración del Modelo de Lenguaje}

% TODO: Completar con el LLM seleccionado, justificación e integración.
\textit{[Pendiente.]}

\section*{8.6 Interfaz de Consulta}

% TODO: Completar con la interfaz desarrollada (API o UI).
\textit{[Pendiente.]}

% ════════════════════════════════════════════════════════════════════════════
\chapter*{9. Investigación Tecnológica}
\addcontentsline{toc}{chapter}{9. Investigación Tecnológica}

Esta sección documenta la investigación realizada por cada integrante del equipo sobre los
componentes tecnológicos del sistema.

% ─── 9.1 LLM ────────────────────────────────────────────────────────────────
\section*{9.1 Modelos de Lenguaje}

\subsection*{¿Cómo correr un LLM localmente?}

La forma más práctica para este proyecto es usando \textbf{Ollama}, una herramienta que
permite descargar y ejecutar modelos como LLaMA~3 o Mistral con un solo comando, sin
configuración compleja. Desde Python se consume mediante una librería simple.

Alternativamente, \textbf{llama.cpp} con su wrapper de Python \texttt{llama-cpp-python}
ofrece mayor flexibilidad, aunque requiere una configuración más elaborada.

\subsection*{Requisitos de hardware}

Los modelos vienen en distintos tamaños medidos en parámetros. Para desarrollo sin GPU,
los modelos \textbf{cuantizados} (formato \texttt{.gguf}) corren en CPU con 8--16\,GB de
RAM, aunque con menor velocidad.

\begin{table}[H]
  \centering
  \begin{tabular}{lcc}
    \toprule
    \textbf{Modelo} & \textbf{RAM mínima} & \textbf{GPU recomendada} \\
    \midrule
    LLaMA 3 8B (cuantizado) & 8\,GB  & No obligatoria \\
    LLaMA 3 70B             & 40+\,GB & Casi obligatoria \\
    Mistral 7B (cuantizado) & 6\,GB  & No obligatoria \\
    \bottomrule
  \end{tabular}
  \caption{Requisitos de hardware por modelo}
\end{table}

\subsection*{Estrategia de desarrollo vs.\ producción}

\begin{itemize}
  \item \textbf{Durante desarrollo y pruebas:} se recomienda usar una API externa
    (OpenAI, Anthropic) para avanzar en el pipeline sin necesidad de hardware dedicado.
  \item \textbf{En la entrega final:} se utiliza el modelo local (LLaMA~3 o Mistral)
    para no depender de servicios externos.
\end{itemize}

Otra opción liviana es \textbf{Phi-3 mini} de Microsoft, que corre en hardware modesto
y tiene buen rendimiento para tareas de Q\&A sobre documentos.

% ─── 9.2 Base de datos vectorial ────────────────────────────────────────────
\section*{9.2 Base de Datos Vectorial}

\subsection*{Comparación de herramientas}

\begin{table}[H]
  \centering
  \small
  \begin{tabular}{lcccc}
    \toprule
    \textbf{Característica} & \textbf{ChromaDB} & \textbf{FAISS} & \textbf{Weaviate} & \textbf{Qdrant} \\
    \midrule
    Tipo              & BD Vectorial   & Librería índice & BD + GraphQL & BD moderna \\
    Persistencia      & Sí (SQLite)    & Manual          & Sí           & Sí \\
    Servidor requerido & No            & No              & Sí (Docker)  & Sí (Docker) \\
    Metadatos/filtros & Sí            & No nativo       & Sí (avanzado)& Sí (potente) \\
    Escalabilidad     & Media          & Alta            & Alta         & Alta \\
    Licencia          & Apache 2.0     & MIT             & BSD-3        & Apache 2.0 \\
    Ideal para        & Prototipado    & Investigación   & Enterprise   & Producción \\
    \bottomrule
  \end{tabular}
  \caption{Comparación de bases de datos vectoriales}
\end{table}

\subsection*{Recomendación para el proyecto}

Dado que los documentos del sistema tienen metadatos estructurados (\texttt{TIPO},
\texttt{ÁREA}, \texttt{AÑO}, \texttt{POZO}, \texttt{SECCIÓN}), se pueden realizar
consultas filtradas del tipo:

\begin{itemize}
  \item \textit{``¿Qué pozos tuvieron problemas de compartimentación?''} →
    \texttt{TIPO = Caracterización Estructural}
  \item \textit{``¿Qué pasó con los inyectores en 2024?''} →
    \texttt{AÑO = 2024 + POZO contiene PI-}
  \item \textit{``Todos los informes de seguimiento de inyección''} →
    \texttt{TIPO = Seguimiento de Inyección}
\end{itemize}

Para \textbf{prototipado local} se recomienda \textbf{ChromaDB} (modo embebido, sin
Docker). Para \textbf{producción con filtros complejos sobre metadatos}, \textbf{Qdrant}
es la opción más adecuada dado el nivel de estructura de los documentos del dominio.

\begin{lstlisting}[caption={Ejemplo de uso con ChromaDB}]
import chromadb
from chromadb.utils import embedding_functions

client = chromadb.PersistentClient(path="./db")
ef = embedding_functions.SentenceTransformerEmbeddingFunction(
    "all-MiniLM-L6-v2"
)
collection = client.get_or_create_collection(
    "documentos", embedding_function=ef
)

# Indexar
collection.add(
    documents=["Texto del documento..."],
    metadatas=[{"tipo": "EWR", "pozo": "PM-104", "anio": 2012}],
    ids=["doc1"]
)

# Consultar con filtro
resultados = collection.query(
    query_texts=["perdida de circulacion en Quintuco"],
    n_results=3,
    where={"tipo": "EWR"}
)
\end{lstlisting}

% ─── 9.3 Embeddings ─────────────────────────────────────────────────────────
\section*{9.3 Embeddings --- Sol}

% TODO: Sol completa esta sección con su investigación sobre embeddings.
\textit{[Pendiente --- Sol Urbano completará esta sección.]}

% ─── 9.4 ETL y OCR ──────────────────────────────────────────────────────────
\section*{9.4 ETL y OCR --- Tomás}

% TODO: Tomi completa esta sección con su investigación sobre ETL y OCR.
\textit{[Pendiente --- Tomás Lanfranco completará esta sección.]}

% ════════════════════════════════════════════════════════════════════════════
\chapter*{10. Riesgos Materializados y Lecciones Aprendidas}
\addcontentsline{toc}{chapter}{10. Riesgos Materializados y Lecciones Aprendidas}

% TODO: Completar durante y al finalizar el desarrollo.
\textit{[Pendiente.]}

% ════════════════════════════════════════════════════════════════════════════
\chapter*{11. Metodología Aplicada}
\addcontentsline{toc}{chapter}{11. Metodología Aplicada}

El proyecto se desarrolló bajo una metodología incremental, organizada en iteraciones que
incluyeron actividades de análisis, implementación, pruebas y mejora continua. Se utilizó
Scrum para la gestión del trabajo, organizando el desarrollo en sprints de una a dos semanas,
con un tablero Kanban en GitHub Projects para el seguimiento de tareas.

% TODO: Completar con reflexión sobre cómo se aplicó la metodología en la práctica.
\textit{[Completar con experiencia real durante el desarrollo.]}

% ════════════════════════════════════════════════════════════════════════════
\chapter*{12. Experimentación y/o Validación}
\addcontentsline{toc}{chapter}{12. Experimentación y/o Validación}

\section*{12.1 Pruebas Unitarias}

% TODO: Completar con resultados de pruebas unitarias por módulo.
\textit{[Pendiente.]}

\section*{12.2 Pruebas de Integración}

% TODO: Completar con resultados de pruebas de integración.
\textit{[Pendiente.]}

\section*{12.3 Pruebas Funcionales End-to-End}

% TODO: Completar con consultas de prueba y resultados obtenidos.
\textit{[Pendiente.]}

\section*{12.4 Métricas de Calidad}

% TODO: Completar con resultados de Top-K accuracy, latencia, tasa de error OCR.
\textit{[Pendiente.]}

% ════════════════════════════════════════════════════════════════════════════
\chapter*{13. Cronograma de Actividades Realizadas}
\addcontentsline{toc}{chapter}{13. Cronograma de Actividades Realizadas}

% TODO: Completar con el cronograma real vs. planificado.
\textit{[Pendiente --- comparar con el plan de actividades del anteproyecto.]}

% ════════════════════════════════════════════════════════════════════════════
\chapter*{14. Impactos Sociales y Ambientales del Proyecto}
\addcontentsline{toc}{chapter}{14. Impactos Sociales y Ambientales del Proyecto}

\section*{14.1 Impacto Económico}

El sistema propuesto permite optimizar el uso de la información disponible, reduciendo la
repetición de errores operativos y mejorando la toma de decisiones en la operación de
cuencas maduras.

\section*{14.2 Impacto Social}

Facilita el acceso al conocimiento técnico acumulado, promoviendo una mejor transferencia
de conocimiento experto dentro del ámbito profesional y reduciendo la dependencia de
personas específicas como repositorios de conocimiento tácito.

\section*{14.3 Impacto Ambiental}

La mejora en la toma de decisiones operativas contribuye a reducir errores que puedan
generar efectos ambientales negativos, favoreciendo un uso más eficiente de los recursos
en las operaciones de extracción.

% ════════════════════════════════════════════════════════════════════════════
\chapter*{15. Trabajos Futuros}
\addcontentsline{toc}{chapter}{15. Trabajos Futuros}

% TODO: Completar con posibles extensiones y mejoras identificadas durante el desarrollo.
\textit{[Pendiente.]}

% ════════════════════════════════════════════════════════════════════════════
\chapter*{16. Conclusiones}
\addcontentsline{toc}{chapter}{16. Conclusiones}

% TODO: Completar al finalizar el proyecto.
\textit{[Pendiente.]}

% ════════════════════════════════════════════════════════════════════════════
\chapter*{Anexos}
\addcontentsline{toc}{chapter}{Anexos}

\section*{A. Estructura del Repositorio}

\begin{lstlisting}[language=bash, caption={Estructura del proyecto}]
rag-knowledge-system/
├── data/
│   ├── raw/            # Documentos originales
│   ├── processed/      # Texto procesado y segmentado
│   └── vectorstore/    # Base de datos vectorial
├── src/
│   ├── etl/            # Ingesta, OCR y chunking
│   ├── embeddings/     # Generacion de vectores
│   ├── retrieval/      # Busqueda semantica
│   ├── llm/            # Integracion del modelo de lenguaje
│   └── interface/      # API o interfaz de consulta
└── tests/              # Pruebas unitarias y de integracion
\end{lstlisting}

\section*{B. Tipos de Documentos Procesados}

\begin{table}[H]
  \centering
  \begin{tabular}{ll}
    \toprule
    \textbf{Tipo} & \textbf{Descripción} \\
    \midrule
    EWR           & Informe Final de Perforación (End of Well Report) \\
    WORKOVER      & Informe Técnico de Workover \\
    ESTIMULACION  & Informe de Estimulaciones y Fracturas \\
    FALLA         & Informe de Análisis de Falla \\
    PARTE\_DIARIO  & Parte Diario de Operaciones \\
    INYECCION     & Informe de Seguimiento de Inyección \\
    GEOLOGIA      & Informe de Caracterización Estructural \\
    PVT           & Estudios PVT y Análisis de Fluidos \\
    EOR           & Informes de Proyectos de Recuperación Mejorada \\
    SEDIMENTOLOGIA& Estudios Sedimentológicos \\
    \bottomrule
  \end{tabular}
  \caption{Tipos de documentos del sistema}
\end{table}

% ─── Bibliografía ────────────────────────────────────────────────────────────
\chapter*{Referencias}
\addcontentsline{toc}{chapter}{Referencias}

\begin{enumerate}
  \item Lewis, P., et al.\ (2020). Retrieval-augmented generation for knowledge-intensive
    NLP tasks. \textit{NeurIPS}. \url{https://arxiv.org/abs/2005.11401}

  \item Dubey, A., et al.\ (2024). The Llama 3 herd of models. \textit{arXiv}.
    \url{https://arxiv.org/abs/2407.21783}

  \item Ma, Z., et al.\ (2024). Information extraction from historical well records using
    a large language model. \textit{Scientific Reports}.
    \url{https://doi.org/10.1038/s41598-024-81846-5}

  \item Guevara, J., et al.\ (2017). A data-driven workflow for predicting horizontal well
    production using vertical well logs. \textit{arXiv}.
    \url{https://arxiv.org/abs/1705.06556}

  \item Sinha, R., \& B.\ S.\ (2025). Digitization of document and information extraction
    using OCR. \url{https://arxiv.org/abs/2506.11156}
\end{enumerate}

\end{document}
